\section{第二部分：非形式化证明——草稿-草图-证明}

\subsection{动机：结合高级推理与低级证明}

\begin{knowledgebox}{Sledgehammer：自动化证明工具}
Isabelle 中的 Sledgehammer 是一种"锤子"（Hammer）工具，它将问题发送给外部自动定理证明器（一阶逻辑、高阶逻辑、SMT 求解器），非常擅长填补证明中的小缺口，但无法处理需要创造性的高级证明策略。
\end{knowledgebox}

\subsection{Draft, Sketch, Prove（DSP）方法}

\begin{importantbox}{三步证明流程}
\begin{enumerate}
    \item \textbf{Draft}（草稿）：LLM 生成非形式化证明（自然语言）
    \item \textbf{Sketch}（草图）：LLM 将非形式化证明翻译为形式化草图，留下待填充的缺口
    \item \textbf{Prove}（证明）：低级证明器（如 Sledgehammer）填补所有缺口
\end{enumerate}
若所有缺口被填补，则获得完整的形式化证明。
\end{importantbox}

搜索策略不同于传统的树搜索：生成多个非形式化草稿和/或多个形式化草图翻译，尝试填补缺口，任一成功即得到证明。

\subsection{实验结果}

在 mini-F2F 上（CodeX/Minerva 时代）：
\begin{itemize}
    \item 随着采样次数增加，性能显著提升
    \item 采样足够多时，LLM 生成的非形式化证明甚至超过了使用人类手写证明的表现
\end{itemize}

